% from ega4-2-fr.tex
\oldpage[IV-2]{34}
générique, $i:X_\eta\to X$ le morphisme canonique. Soient $\mathcal F$ un
$\mathcal O_X$-Module quasi-cohérent, $\mathcal F_\eta=i^*(\mathcal F)$,
$\mathcal G'$ un $\mathcal O_{X_\eta}$-Module quotient de $\mathcal F_\eta$,
et soit $\overline{\mathcal G'}$ le $\mathcal O_X$-Module image de $\mathcal F$
par l'homomorphisme composé $(0_{\mathrm I},4.4.3.2)$
\[
\mathcal F\xrightarrow{\rho_{\mathcal F}}i_*(i^*(\mathcal F))
\longrightarrow i_*(\mathcal G').
\]
Alors $\overline{\mathcal G'}$ est un $\mathcal O_X$-Module quotient de
$\mathcal F$, quasi-cohérent et $f$-plat, tel que
$i^*(\overline{\mathcal G'})=\mathcal G'$, et c'est le seul
$\mathcal O_X$-Module quotient de $\mathcal F$ ayant ces propriétés.
\end{proposition}

\begin{proof}
Comme $i$ est quasi-compact et quasi-séparé (1.1.2 et 1.2.2), il résulte de
(1.7.4) que pour tout $\mathcal O_{X_\eta}$-Module quasi-cohérent
$\mathcal H'$, $i_*(\mathcal H')$ est un $\mathcal O_X$-Module
quasi-cohérent; en outre, pour tout ouvert $U$ de $X$, on a
\[
(i_*(\mathcal H'))|U=(i|U\cap X_\eta)_*(\mathcal H'|U\cap X_\eta)
\]
par définition $(0_{\mathrm I},3.4.1)$. Si l'on démontre la proposition
lorsque $X$ et $Y$ sont affines, elle en résultera par recollement dans le cas
général, vu l'assertion d'unicité valable dans le cas affine. Autrement dit,
on est ramené à prouver le

\begin{lemma}[2.8.1.1]
\label{IV.2.8.1.1-fr}
Soient $A$ un anneau noethérien régulier $(0,17.3.6)$, intègre de dimension
1, $K$ son corps des fractions, $M$ un $A$-module, $N'$ un $K$-module
quotient de $M_{(K)}=M\otimes_A K$ par un sous-$K$-module $P'$,
$\overline{N'}$ l'image de $M$ par l'homomorphisme composé
$M\to M_{(K)}\to N'$. Alors $\overline{N'}$ est un $A$-module plat, et est
l'unique module quotient $N$ de $M$ qui soit un $A$-module plat et tel que le
noyau de l'homomorphisme surjectif $M_{(K)}\to N_{(K)}$ soit égal à $P'$.
\end{lemma}

Comme pour tout idéal maximal $\mathfrak m$ de $A$, $A_{\mathfrak m}$ est un
anneau local régulier de dimension 1, donc un anneau de valuation discrète,
il revient au même de dire qu'un $A$-module $N$ est plat ou qu'il est
\emph{sans torsion} $(0_{\mathrm I},6.3.4)$. Comme $N'$ est un $K$-espace
vectoriel, c'est un $A$-module sans torsion, donc il en est de même de
$\overline{N'}$, sous-module de $N'$; en outre, il est immédiat de vérifier
que $\overline{N'}_{(K)}$ s'identifie à $N'$. Inversement, si $N$ est un
$A$-module quotient de $M$ ayant les propriétés de l'énoncé, le fait que $N$
soit un $A$-module plat entraîne que l'homomorphisme canonique
$N\to N_{(K)}=N\otimes_A K$ est injectif. Comme $N_{(K)}$ s'identifie à
$N'$, la conclusion résulte de la commutativité du diagramme
\[
\xymatrix{
M \ar[r] \ar[d] & N \ar[d] \\
M_{(K)} \ar[r] & N_{(K)}.
}
\]
\end{proof}

\begin{corollary}[2.8.2]
\label{IV.2.8.2-fr}
Sous les conditions de (2.8.1), pour que $\mathcal F$ soit $f$-plat, il faut
et il suffit que l'homomorphisme canonique
\[
\mathcal F\longrightarrow i_*(\mathcal F_\eta)=i_*(i^*(\mathcal F))
\]
soit injectif.
\end{corollary}

\begin{env}[2.8.3]
\label{IV.2.8.3-fr}
La formation du $\mathcal O_X$-Module $\overline{\mathcal G'}$ est
\emph{fonctorielle en $\mathcal F$ et $\mathcal G'$} : de façon précise, si
$\mathcal F_1$, $\mathcal F_2$ sont deux $\mathcal O_X$-Modules
quasi-cohérents, $u:\mathcal F_1\to\mathcal F_2$ un
$\mathcal O_X$-homomorphisme, $\mathcal G'_i$ un
$\mathcal O_{X_\eta}$-Module quotient de $(\mathcal F_i)_\eta$ $(i=1,2)$ et
$v:\mathcal G'_1\to\mathcal G'_2$ un homomorphisme rendant commutatif le
diagramme
\[
\xymatrix{
(\mathcal F_1)_\eta \ar[r]^{i^*(u)} \ar[d] &
(\mathcal F_2)_\eta \ar[d] \\
\mathcal G'_1 \ar[r]_v & \mathcal G'_2
}
\]

